% Nội dung báo cáo được chuyển từ bản Markdown đã đối chiếu.
\chapter{KẾT LUẬN VÀ CẢM NGHĨ CỦA NHÓM}

\section{Kết luận}

Đề tài đã kết hợp nghiên cứu lịch sử với việc xây dựng một sản phẩm học liệu số có thể truy cập và sử dụng trên trình duyệt. Phần nội dung hệ thống hóa sự lãnh đạo của Đảng qua năm giai đoạn kháng chiến chống Mỹ, cứu nước; phần sản phẩm chuyển nội dung đó thành dòng thời gian tương tác, kho tư liệu, bài tổng kết và hoạt động trắc nghiệm có giải thích. Website vì vậy là kết quả ứng dụng trực tiếp của quá trình chọn lọc, tổ chức và trình bày tri thức mà nhóm đã thực hiện.

Kết quả nổi bật của nhóm không chỉ là một địa chỉ web được triển khai, mà còn là sự kết nối giữa các hình thức tiếp cận: đọc để nắm tiến trình, quan sát để hình dung chứng tích và tự kiểm tra để nhận ra nội dung cần củng cố. Qua đó, mục tiêu “từ đường lối đến thực tiễn” được thể hiện cả trong nội dung lịch sử lẫn cách tổ chức hoạt động học tập. Những chức năng đã kiểm tra cho thấy sản phẩm đáp ứng các thao tác cơ bản trong phạm vi đề tài; việc đánh giá hiệu quả học tập thực tế cần được tiếp tục bằng phản hồi và khảo sát người sử dụng.

\section{Cảm nghĩ sau quá trình khảo sát và xây dựng sản phẩm}

Qua việc tìm hiểu chứng tích tại Bảo tàng Chứng tích Chiến tranh, nhóm nhận thức rõ hơn rằng các mốc lịch sử trong giáo trình gắn với đời sống, sự hy sinh và những mất mát cụ thể của con người. Các hiện vật, ảnh tư liệu và nội dung trưng bày giúp việc học không dừng lại ở tên gọi của một chiến lược hay kết quả của một chiến dịch. Đằng sau đó là những nỗ lực bền bỉ của nhân dân, sự phối hợp giữa hậu phương và tiền tuyến, cùng khát vọng độc lập, thống nhất và hòa bình.

Những ghi nhận ấy làm sâu sắc lòng biết ơn của nhóm đối với các thế hệ đi trước, đồng thời nhắc nhở trách nhiệm khi chuyển tư liệu thành nội dung số. Một hình ảnh không nên chỉ được dùng để tạo ấn tượng; nó cần có bối cảnh và chú thích phù hợp. Một câu hỏi trắc nghiệm không chỉ cần đáp án đúng mà còn cần giải thích rõ. Với nhóm, thái độ nghiêm túc trong kiểm chứng và trình bày kiến thức cũng là cách thể hiện sự tôn trọng lịch sử và người học.

Quá trình xây dựng website còn cho thấy giá trị của hợp tác. Nội dung, thiết kế và lập trình chỉ tạo thành sản phẩm khi được phối hợp và rà soát cùng nhau. Các thành viên có cơ hội rèn luyện kỹ năng phân công, tổng hợp tài liệu, trao đổi và chịu trách nhiệm đối với phần việc chung. Việc vận dụng công nghệ vào môn học giúp nhóm nhận thấy chuyên môn của sinh viên có thể góp phần làm cho học liệu dễ tiếp cận hơn, nhưng không thay thế yêu cầu nghiên cứu nghiêm túc.

Từ trải nghiệm đó, nhóm xác định trách nhiệm của người trẻ bắt đầu từ những việc cụ thể: học tập có mục tiêu, sử dụng thông tin có kiểm chứng, tôn trọng con người và đóng góp bằng năng lực của mình. Trân trọng hòa bình không chỉ là cảm xúc trước chứng tích chiến tranh mà còn là ý thức gìn giữ, xây dựng cuộc sống hôm nay và ứng xử có trách nhiệm trong cộng đồng.

\section{Lời cảm ơn}

Nhóm xin trân trọng cảm ơn Ban Giám hiệu Trường Đại học Khoa học Tự nhiên, Đại học Quốc gia Thành phố Hồ Chí Minh; quý Thầy, Cô phụ trách học phần Lịch sử Đảng Cộng sản Việt Nam đã tạo điều kiện để sinh viên kết hợp học tập lý luận với tìm hiểu thực tế. Nhóm cảm ơn Bảo tàng Chứng tích Chiến tranh đã gìn giữ, trưng bày và giới thiệu hệ thống tư liệu giúp công chúng tiếp cận lịch sử từ những chứng tích cụ thể.

Nhóm cũng cảm ơn sự phối hợp và đóng góp của các thành viên trong quá trình biên soạn nội dung, khảo sát, xây dựng và hoàn thiện website. Dù đã nỗ lực thực hiện, sản phẩm và báo cáo vẫn có thể còn hạn chế. Nhóm mong nhận được góp ý của giảng viên và người học để tiếp tục hiệu đính, hoàn thiện học liệu theo hướng chính xác, rõ ràng và hữu ích hơn.

\clearpage
